\section{Results and Discussion}
\label{sec:results}
No research run has yet been verified for inclusion. Historical scores and engineering fixtures are excluded. The final comparison will report each arm's evaluated population, attempted and valid counts, and available micro and macro accuracy. Generation scores will remain unavailable until judgments under an approved protocol exist.

\ifpaperresults
\begin{table*}[t]
\centering\small
\PaperResultsTable
\caption{Selection scores and execution coverage from the explicitly selected preliminary development run. Free generation remains unjudged; unavailable values are not zero.}
\label{tab:run-results}
\end{table*}
\begin{figure*}[t]
\centering
\PaperSelectionFigure
\caption{Preliminary development selection accuracy for run \PaperRunID. Micro divides by all requested items; macro weights types equally. Partial arms are labeled, and failed requests remain in the denominator.}
\label{fig:selection-results}
\end{figure*}
\else
\begin{table}[t]
\centering\small
\begin{tabular}{@{}lp{0.58\columnwidth}@{}}
\toprule
Planned display & Question it will address \\
\midrule
Main results & What is each arm's accuracy and completion status on the same items? \\
Selection plot & How do micro and macro accuracy differ across selection systems? \\
Paired analysis & Which occurrences change correctness between generation and selection? \\
\bottomrule
\end{tabular}
\caption{Reserved results structure. No numerical entries are inferred from pending runs.}
\label{tab:planned-results}
\end{table}
\fi

The discussion will distinguish semantic mistakes from output-format and service failures, relate disagreements to inventory coverage, and examine whether aggregate performance is concentrated in a few types. It will interpret the paired contrasts in light of the response-format difference and the limited sense diversity of development data. There is presently no empirical conclusion about the preferred formulation or relative system performance.
